\MajorSection{القسم التاسع}

\Couplet{٢٢٤}{فكيف ينظم الإنسان حياته، حتى إذا انتهى تعداد سنيه؛}{مضى في طريقه كضيف شبع؟ وكيف يحافظ على مساره المستوي؟}{}

\Couplet{٢٢٥}{رغم الكتابة التي تحتفظ بها الجمجمة، ورغم اللوح والقلم؛}{ورغم القدر الذي يلعب بنا ليطرحنا، ورقعته العالم وقطعه البشر؟}{\BurtonNote{حاشية برتون: رموز القسمة أو المصير.}}

\Couplet{٢٢٦}{وكيف، حين يخفت نور الحياة ووهجها في ظلمة تتكاثف؛}{يسخر الفاني من شوكة الموت، ويحتقر ظفر القبر؟}{}

\Couplet{٢٢٧}{طريق واحد ومسلكان، ونهاية واحدة هي القبر؛ هذا يمر في السهل المزهر؛}{وذاك يواجه الشجيرات والمنحدر والصخر، في الشمس والريح والثلج والمطر.}{}

\Couplet{٢٢٨}{ومن يسلك الأول عليه أن ينظر إلى أسفل، وأن يعد حياته الكل في الكل؛}{وألا يرى أعالي يمكن للإنسان أن يرتفع إليها، ولا أعماقًا يمكن أن يسقط فيها.}{}

\Couplet{٢٢٩}{وعليه أن يرى الله في صورة آدم، وأن يعبد الصانع في المصنوع؛}{راضيًا بالاستلقاء في ابتسامة مايا، في مسرّات الألم، وفي أنوار الظل.}{\BurtonNote{حاشية برتون: مايا هي الوهم.}}

\Couplet{٢٣٠}{يخرق الشريعة ويحرق الكتاب، ويرسل الملّا من جديد إلى المدرسة؛}{ويضحك من لحى الرجال القديسين، ويسمّي النبي غبيًا وأحمق.}{}

\Couplet{٢٣١}{ويعانق خصر السرو المستدق، ويبرّد قدميه على صدر الجدول المتموج؛}{ويبتسم في عيني النرجس المكلومتين بالحب، ويتمتع برقصة النرجس الأصفر.}{}

\Couplet{٢٣٢}{ويذوب في ضوء الفجر الزعفراني ليسمع أنين الحمامة؛}{ويتلذذ بألوان الغروب الأرجوانية حين يخطب البلبل حب الوردة.}{}

\Couplet{٢٣٣}{ويجد المرح والفرح في كأس جمشيد، ويداعب بنت الكرمة؛}{ويأمر غلام السقاية الجميل أن يقول: «سيدي، أحمل إليك خمرًا ياقوتية!»}{\BurtonNote{حاشية برتون: لكي تتمتع الحواس كلها، حتى الأذن.}}

\Couplet{٢٣٤}{ويرتشف الندى من شفتي الفتاة، ويمسح النضارة عن جبين العذراء؛}{تلك هي سعادته الجسدية، وهو يسعى إلى معرفة الصانع بواسطة المصنوع.}{}

\Couplet{٢٣٥}{جرّبتها كلها، فوجدتها كلها متشابهة فاترة، كئيبة جافة؛}{وتتقلب معدتي عند التفكير بها، فأخاطب نفسي وأصيح:}{}

\Couplet{٢٣٦}{خير منها الكدود والآلام التي لا تُحصى، والتي تجعل الرجل وفيًا لرجولته؛}{لتكن هذه القاعدة التي تهدي الحياة، وهذه القوانين لي ولك.}{}

\Couplet{٢٣٧}{شنّ على الجهل حربًا أبدية، واجهد دائمًا في معرفة نفسك؛}{فجهلك بجهلك أشرس أعدائك، وأشد أسباب هلاكك.}{}

\Couplet{٢٣٨}{إنه يثلم إحساسك ويبلّد ذوقك، ويصم أذنيك ويعمي عينيك؛}{ويخلق الشيء الذي لم يكن قط، ويتحدى الشيء الذي يكون دائمًا.}{}

\Couplet{٢٣٩}{وتلك الذرة المتناهية اللامتناهية، التي تشكّل النقطة الوسطى في دائرتك؛}{والمكتفية كل الاكتفاء بنفسها، وغير الموجودة للذوات الأخرى؛}{}

\Couplet{٢٤٠}{تجد العالم قويًا، بمقدار ما هي صغيرة؛ ومع ذلك لا بد من خوض المعركة غير المتكافئة؛}{هنا عمالقة لا يُحصَون، وهناك قبضة غبار وكتلة طين.}{}

\Couplet{٢٤١}{نعم، رغم أحلامك كلها بالسلام، لا بد من خوض الصراع غير العادل؛}{حيث يمكنك أن تتعلم أنبل معرفة: أن تعرف أن كل ما نعرفه لا شيء.}{}

\Couplet{٢٤٢}{كن وفيًا لطبيعتك ولنفسك؛ فلا ترجُ الصيت ولا تخشَه، ولا ترجُ سوء الصيت ولا تخشَه؛}{يكفيك ذلك الصوت الخافت الساكن الذي يرعد دائمًا في أذنك الباطنة.}{}

\Couplet{٢٤٣}{اطلب الاستحسان من رضاك عن نفسك؛ فما لا يعرفه الناس تعرفه أنت؛}{وازدرِ كل صنم يقيمه الآخرون، وانحنِ أمام مثالك أنت.}{}

\Couplet{٢٤٤}{كن إله نفسك؛ واجعل الذات حرة سمحة كالهواء المحيط؛}{وليكن فكرك لك إمبراطورية؛ وحطّم كل قفل وقضيب يحبسها.}{}

\Couplet{٢٤٥}{أدِّ ما تدين به دائمًا لنفسك من واجب؛ هنا تلتقي الواجبات كلها وتمتزج؛}{بأوسع معانيها، من غير اهتمام بما بدأ، أو بما سينتهي.}{}

\Couplet{٢٤٦}{وهكذا، حين تتأمل الصور الشبحية التي كانت لك في الماضي الضبابي؛}{يمكنك، بفخر صادق، أن ترفض أن تصير من جديد ما كنت عليه.}{}

\Couplet{٢٤٧}{وحين تمد بصرك في صف السنين، لا تخشَ أن ترى ذاتك المقبلة؛}{راضيًا بالحياة، راضيًا بالموت، كما لو أن الاختيار لا شأن له عندك.}{}

\Couplet{٢٤٨}{لا تغذّ فكرك بالفكر نفسه؛ ولا تستدر من الشمس والنور لتحدّق}{في أروقة معتمة رُصفت بالقبور، حيث تتعفن عظام الأيام الغابرة.}{}

\Couplet{٢٤٩}{قال الحكماء: «لا تأكل قلبك، ولا تنح على الماضي، الماضي المدفون»؛}{افعل ما تفعل، وكن قويًا شجاعًا؛ ومثل النجم، لا تسترح ولا تتعجل.}{}

\Couplet{٢٥٠}{انتزع العجوز من صدرك؛ وكن جلدًا في الشقاء، وصلبًا في الهناء؛}{وافعل الخير لأن فعل الخير خير؛ وازدرِ رشوة الجنة وتهديد الجحيم.}{}

\Couplet{٢٥١}{طلب الحق وإسعاد القلب: ذلك هو الناموس الأعلى للحياة؛}{وبه يتفاوت قدر الرجل، بين رجل من ذهب ورجل من قش.}{}

\Couplet{٢٥٢}{لا تنظر إلى ذلك الشيء في البشر الذي يثير الكراهية أو الاحتقار أو الخصام؛}{لَدودة عزرائيل خير من موت يمشي في صورة حياة.}{\BurtonNote{حاشية برتون: عزرائيل ملاك الموت.}}

\Couplet{٢٥٣}{تأمل بني جنسك بوصفهم واحدًا، تتحد حاجاته في الكل الإنساني العظيم؛}{الإنسان الذي يرتفع من الأرض عاليًا ليطلب سماوات الحياة في النور.}{\BurtonNote{حاشية برتون: «الإنسان العظيم» عند الأخنوخيين والمورمون.}}

\Couplet{٢٥٤}{واعتبر الإنسانية إنسانًا واحدًا، لا يزال عذابه الشامل}{يكد ويجاهد ليبلغ الغاية التي يكف فيها العذاب عن الوجود.}{}

\Couplet{٢٥٥}{آمن بالأشياء كلها، ولا تؤمن بشيء؛ ولا تحكم، ولا تحرّف الفكر بـ«الوقائع»؛}{وانظر بوضوح واسمع بوضوح، وإن بدت الحياة مايا وسرابًا وحلمًا وعدمًا.}{}

\Couplet{٢٥٦}{انصرف عن «لماذا» واطلب «كيف»؛ فالإله والآلهة المتربعة في الأعالي؛}{صامتة جميعًا، ولا تزال صامتة؛ لا تسمع صوتك، ولا تتفضل بجواب.}{}

\Couplet{٢٥٧}{الآن، تلك النقطة التي لا تتجزأ، والمنثورة على طول خط لا متناهٍ؛}{لا تقع نهايتاه في مكان، هي كل ما لك، ذلك الكل الضئيل الذي تسمّيه لك.}{}

\Couplet{٢٥٨}{لعل للناموس واهبًا؛ دع الأمر، دع الأمر! ماذا يمكنك أن تعرف؛}{لقد جاءت شعوب لا تُحصى ومضت، وقد رآها أبو الهول هذا تأتي وتمضي.}{}

\Couplet{٢٥٩}{ولعل الناموس الذي يحكم العالم يتيح للإنسان أوسع مجال؛}{ولعل القدر كلمة عند المؤمن بإله شخصي، خاضعة لمصادفة الإنسان وتغيّره.}{}

\Couplet{٢٦٠}{وقد تجد «أنا» هذه حياة مقبلة، نسخة أنبل من حياتنا؛}{حيث يحل كل لغز، وتُعرف كل معرفة.}{}

\Couplet{٢٦١}{حيث يكون للإنسان أن يرى كاملًا ما يراه على الأرض جزءًا؛}{وحيث لا يثقل التغيّر الفكر، ولا يجرح الرجاء المؤجّل القلب.}{}

\Couplet{٢٦٢}{ولكن الزهرة الذابلة والورقة الساقطة لن تزيّنا الشجرة الأم مرة أخرى؛}{والإنسان، إذا أسقطته شجرة الحياة مرة، فأي رجاء له في حياة أخرى؟}{}

\Couplet{٢٦٣}{يمكن إصلاح الوعاء المحطم، ويمكن للعود المشقوق أن يصوّت من جديد؛}{لكن من يصلح طين الإنسان، ومن يعيد إليه النفَس المسلوب؟}{}

\Couplet{٢٦٤}{ستدق الساعة المحطمة من جديد، وسيزمر القصب المكسور مرة أخرى؛}{أما نحن فنموت، والموت واحد: قضاء البهائم وقضاء البشر.}{}

\Couplet{٢٦٥}{فإن ختمت النيرفانا حياتنا بالعدم، فلعل هذا هو الأفضل؛}{لقد نالت أخيرًا كدودك ومتاعبك، وعوزك وشقاؤك، جزاءها: الراحة.}{\BurtonNote{حاشية برتون: الفناء النسبي.}}

\Couplet{٢٦٦}{كفّ، يا عبدو، كفّ! لقد غُنّيت أغنيتك، ولا تظن أن الكسب جائزة المغني؛}{إلى أن يعدّ الرجال الجهل خطيئة مهلكة، وإلى أن يستحق الإنسان لقبه: العاقل.}{\BurtonNote{حاشية برتون: \textenglish{Homo sapiens}، الإنسان العاقل.}}

\Couplet{٢٦٧}{في أيام مقبلة، أيام بطيئة الانبثاق، حين تتفضل الحكمة بالإقامة بين البشر؛}{قد توقظ هذه الأصداء لصوت سكت منذ زمن طويل نغمة تجيبها.}{}

\Couplet{٢٦٨}{امضِ الآن في طريقك، بجبين هادئ، ولا تخشَ أن تروي حكايتك المتواضعة؛}{همسات ريح الصحراء، ورنين جرس الجمل.}{\BurtonNote{ختام النص بالعبرية: سلام، مكتوبة بالحروف \textenglish{ShLM}.}}
